\section{Case-Study 2: Ergebnisse der FMEA innerhalb der Risk Gates}

\subsection{Risk Gate 2: Kriterien zur Bildung der RPZ}
\begin{table}[H]
\centering
\renewcommand{\arraystretch}{1.3}
\begin{tabular}{|p{3cm}|p{4.9cm}|p{3cm}|p{4.9cm}|}
\hline
\textbf{Risiko} & \textbf{Auftritts\-wahrscheinlichkeit A} & \textbf{Bedeutung B} & \textbf{Entdeckungs\-wahrscheinlichkeit E} \\
\hline

Mangelnde Datenverfügbarkeit und Zugriffshürden 
& 5 (gelegentliches Auftreten)
& 5 (Störung im Prozess)
& 2 (zwangsläufige Entdeckung) \\
\hline

Bias 
& 3 (gering)
& 3 (unbedeutend)
& 2 (zwangsläufige Entdeckung im folgenden Prozess) \\
\hline

Rechtliche Verstöße
& 4 (gelegentliches Auftreten)
& 10 (Verletzung von Vorschriften)
& 9 (sachverständiger Stakeholder wird Fehler entdecken) \\
\hline

Mangelnde Interpretierbarkeit
& 5 (gelegentliches Auftreten)
& 5 (Störung im Prozess)
& 2 (zwangsläufige Entdeckung) \\
\hline

Res\-sour\-cen\-über\-for\-de\-rung
& 3 (gering)
& 3 (unbedeutend)
& 2 (zwangsläufige Entdeckung im folgenden Prozess) \\
\hline

Fehlende Integration der Business Perspektive
& 4 (gelegentliches Auftreten)
& 10 (Verletzung von Vorschriften)
& 9 (sachverständiger Stakeholder wird Fehler entdecken) \\
\hline

\end{tabular}
\caption{Risikoregister: Kriterien zur Bildung der RPZ innerhalb des 2. Risk Gates, eigene Darstellung nach \cite{Bundesministerium_2018}}
\label{tab:cs2_app_rg2_kriterien_rpz}
\end{table}

\subsection{Risk Gate 2: Ergebnisse}

\begin{table}[H]
\centering
\begin{tabular}{|p{9cm}|c|}
\hline
\textbf{Risiko} & \textbf{RPZ} \\ \hline
Mangelnde Datenverfügbarkeit und Zugriffshürden & $(5 \times 5 \times 2) = 50$ \\ \hline
Bias & $(3 \times 3 \times 2) = 18$ \\ \hline
Rechtliche Verstöße & $(4 \times 10 \times 9) = 360$ \\ \hline
Mangelnde Interpretierbarkeit & $(3 \times 5 \times 2) = 50$ \\ \hline
Ressourcenüberforderung & $(2 \times 5 \times 2) = 18$ \\ \hline
Fehlende Integration der Business Perspektive & $(2 \times 3 \times 2) = 360$ \\ \hline
\end{tabular}
\caption{Beispielhafte Risikopriorisierung mittels RPZ, Risk Gate 2, eigene Darstellung nach \cite{Bundesministerium_2018}}
\label{tab:cs2_app_rg2_rpz_risiken}
\end{table}

\subsection{Risk Gate 3: Kriterien zur Bildung der RPZ}

\begin{table}[H]
\centering
\renewcommand{\arraystretch}{1.3}
\begin{tabular}{|p{3cm}|p{4.9cm}|p{3cm}|p{4.9cm}|}
\hline
\textbf{Risiko} & \textbf{Auftritts\-wahrscheinlichkeit A} & \textbf{Bedeutung B} & \textbf{Entdeckungs\-wahrscheinlichkeit E} \\
\hline

Overfitting 
& 6 (gelegentliches Auftreten)
& 7 (eingeschränkte Dienstleistung)
& 3 (Entdeckung des Fehlers mit hoher Wahrscheinlichkeit) \\
\hline

Mangelnde Datenverfügbarkeit und Zugriffshürden
& 3 (geringe Wahrscheinlichkeit)
& 4 (Störung im Prozess)
& 10 (Entdeckung nicht sofort möglich) \\
\hline

Bias
& 3 (gering)
& 3 (unbedeutend)
& 2 (zwangsläufige Entdeckung im folgenden Prozess) \\
\hline

Rechtliche Verstöße
& 3 (geringes Risiko)
& 10 (Verletzung von Vorschriften)
& 9 (sachverständiger Stakeholder wird Fehler entdecken) \\
\hline

Mangelnde Interpretierbarkeit
& 3 (gering)
& 5 (Störung im Prozess)
& 2 (zwangsläufige Entdeckung) \\
\hline

Res\-sour\-cen\-über\-for\-de\-rung
& 2 (gering)
& 3 (unbedeutend)
& 2 (zwangsläufige Entdeckung im folgenden Prozess) \\
\hline

Fehlende Integration der Business Perspektive
& 1 (fast auszuschließen)
& 2 (unbedeutend)
& 2 (sachverständiger Stakeholder wird Fehler entdecken) \\
\hline

\end{tabular}
\caption{Risikoregister: Kriterien zur Bildung der RPZ innerhalb des 3. Risk Gates, eigene Darstellung nach \cite{Bundesministerium_2018}}
\label{tab:cs2_app_rg3_kriterien_rpz}
\end{table}

\subsection{Risk Gate 3: Ergebnisse}

\begin{table}[H]
\centering
\begin{tabular}{|p{9cm}|c|}
\hline
\textbf{Risiko} & \textbf{RPZ} \\ \hline
Overfitting & $(6 \times 7 \times 3) = 126$ \\ \hline
Mangelnde Datenverfügbarkeit und Zugriffshürden & $(3 \times 4 \times 10) = 120$ \\ \hline
Bias & $(2 \times 3 \times 2) = 12$ \\ \hline
Rechtliche Verstöße & $(3 \times 10 \times 9) = 270$ \\ \hline
Mangelnde Interpretierbarkeit & $(3 \times 5 \times 2) = 30$ \\ \hline
Ressourcenüberforderung & $(2 \times 3 \times 2) = 12$ \\ \hline
Fehlende Integration der Business Perspektive & $(1 \times 2 \times 2) = 4$ \\ \hline
\end{tabular}
\caption{Beispielhafte Risikopriorisierung mittels RPZ, Risk Gate 3, eigene Darstellung nach \cite{Bundesministerium_2018}}
\label{tab:cs2_app_rg3_rpz_risiken}
\end{table}

\subsection{Risk Gate 4: Kriterien zur Bildung der RPZ}

\begin{table}[H]
\centering
\renewcommand{\arraystretch}{1.3}
\begin{tabular}{|p{3cm}|p{4.9cm}|p{3cm}|p{4.9cm}|}
\hline
\textbf{Risiko} & \textbf{Auftritts\-wahrscheinlichkeit A} & \textbf{Bedeutung B} & \textbf{Entdeckungs\-wahrscheinlichkeit E} \\
\hline

Inkompatibilität der IT-Infrastruktur mit dem Modell
& 3 (gering)
& 7 (eingeschränkte Dienstleistung)
& 1 (zwangsläufige Entdeckung) \\
\hline

Mangelnde Akzeptanz des Endanwenders
& 7 (häufiges Auftreten)
& 5 (Störung im Prozess, Probleme bei manchen Kunden)
& 10 (Entdeckung nicht sofort möglich) \\
\hline

Unklare Vorgehensweise im Fehlerfall
& 7 (häufiges Auftreten)
& 4 (Störung im Prozess, Probleme bei manchen Kunden)
& 10 (Entdeckung nicht sofort möglich) \\
\hline

Overfitting 
& 6 (gelegentliches Auftreten)
& 7 (eingeschränkte Dienstleistung)
& 3 (Entdeckung des Fehlers mit hoher Wahrscheinlichkeit) \\
\hline

Mangelnde Datenverfügbarkeit und Zugriffshürden
& 3 (geringe Wahrscheinlichkeit)
& 4 (Störung im Prozess)
& 10 (Entdeckung nicht sofort möglich) \\
\hline

Bias 
& 3 (gering)
& 3 (unbedeutend)
& 2 (zwangsläufige Entdeckung im folgenden Prozess) \\
\hline

Rechtliche Verstöße
& 3 (geringes Risiko)
& 10 (Verletzung von Vorschriften)
& 9 (sachverständiger Stakeholder wird Fehler entdecken) \\
\hline

Mangelnde Interpretierbarkeit
& 3 (gering)
& 5 (Störung im Prozess)
& 2 (zwangsläufige Entdeckung) \\
\hline

Res\-sour\-cen\-über\-for\-de\-rung
& 2 (gering)
& 3 (unbedeutend)
& 2 (zwangsläufige Entdeckung im folgenden Prozess) \\
\hline

Fehlende Integration der Business Perspektive
& 1 (fast auszuschließen)
& 2 (unbedeutend)
& 2 (sachverständiger Stakeholder wird Fehler entdecken) \\
\hline

\end{tabular}
\caption{Kriterien zur Bildung der RPZ innerhalb des 4. Risk Gates, eigene Darstellung nach \cite{Bundesministerium_2018}}
\label{tab:cs2_app_rg4_kriterien_rpz}
\end{table}

\subsection{Risk Gate 4: Ergebnisse}

\begin{table}[H]
\centering
\begin{tabular}{|p{9cm}|c|}
\hline
\textbf{Risiko} & \textbf{RPZ} \\ \hline
Inkompatibilität mit der IT-Infrastruktur & $(3 \times 7 \times 1) = 21$ \\ \hline
Mangelnde Akzeptanz beim Endanwender & $(7 \times 5 \times 10) = 350$ \\ \hline
Unklare Vorgehensweise im Fehlerfall & $(7 \times 4 \times 10) = 280$ \\ \hline
Mangelnde Datenverfügbarkeit und Zugriffshürden & $(5 \times 5 \times 2) = 50$ \\ \hline
Bias & $(3 \times 3 \times 2) = 18$ \\ \hline
Rechtliche Verstöße & $(3 \times 10 \times 9) = 270$ \\ \hline
\end{tabular}
\caption{Beispielhafte Risikopriorisierung mittels RPZ, Risk Gate 4, eigene Darstellung nach \cite{Bundesministerium_2018}}
\label{tab:cs2_app_rg4_rpz_risiken}
\end{table}
